% Reviewed translation: PDF pages 265--271 / printed pages 251--257.
% The original scan is authoritative; mathematical tokens were checked against it.
\MathBookSourcePage{265}
\subsection{Hellan--Herrmann--Johnson 格式的补充结果}
\label{subsec:incompressible-hhj-additional-results}

本小节把前面的部分结果推广到右端项
$\mathbf f\in L^r(\Omega)^2$, $1<r<2$ 的情形.
问题 \eqref{eq:incompressible-hhj-discrete-problem} 的空间是有限维的,
特别地 $\Phi_h\subset W^{1,\infty}(\Omega)$,
所以当 $\mathbf f$ 仅属于 $L^r(\Omega)^2$ 时该问题仍有唯一解.
因而只需调整连续问题 \eqref{eq:incompressible-hhj-formulation}:
降低张量检验函数的正则性, 同时提高流函数检验函数的正则性.

若 Stokes 问题的解 $(\mathbf u=\operatorname{curl}\psi,p)$ 满足
$\psi\in W^{3,r}(\Omega)$, 则
$\boldsymbol\sigma=(\partial^2\psi/\partial x_i\partial x_j)_{i,j}
\in W^{1,r}(\Omega)^4$.
由 Sobolev 嵌入定理和 trace 定理,
\[
\boldsymbol\sigma\in L^2(\Omega)^4,
\qquad M_n(\boldsymbol\sigma)\in L^{r/(2-r)}(\Gamma_h).
\]
因此把张量空间替换为
\[
\widetilde{\mathcal X}
=\{\boldsymbol\tau\in L^2(\Omega)^4;
\ \boldsymbol\tau|_\kappa\in W^{1,r}(\kappa)^4,
\ \tau_{12}=\tau_{21},
\ M_n(\boldsymbol\tau)\text{ 在 }\Gamma_h\text{ 每条边上连续}\}.
\]
又因 $\mathbf f\in L^r(\Omega)^2$, 把流函数空间替换为
$\widetilde\Psi\cap W^{1,s}(\Omega)$, 其中 $1/s+1/r=1$.
于是 $(\psi,\boldsymbol\sigma)$ 是下列问题的唯一解:
\begin{equation}\label{eq:incompressible-hhj-low-regularity-formulation}
\left\{
\begin{aligned}
b_h(\boldsymbol\sigma,\phi)
&=-\frac1\nu\int_\Omega\mathbf f\cdot\operatorname{curl}\phi\,dx
&&\forall\phi\in\widetilde\Psi\cap W^{1,s}(\Omega),\\
a_h(\boldsymbol\sigma,\boldsymbol\tau)
+b_h(\boldsymbol\tau,\psi)&=0
&&\forall\boldsymbol\tau\in\widetilde{\mathcal X}.
\end{aligned}
\right.
\end{equation}

Theorem~\ref{thm:incompressible-hhj-error} 及其 Corollary 的证明表明,
只需验证相应的正交关系、稳定估计和误差估计仍然成立.
这些结论来自原有误差方程以及
\[
b_h(\boldsymbol\sigma-\Pi_h\boldsymbol\sigma,\phi_h)=0
\quad\forall\phi_h\in\Phi_h,
\qquad
b_h(\boldsymbol\tau_h,\psi-I_h\psi)=0
\quad\forall\boldsymbol\tau_h\in\mathcal X_h.
\]
当 $\psi\in W^{3,r}(\Omega)$ 且
$\boldsymbol\sigma\in W^{1,r}(\Omega)^4$ 时,
$I_h\psi$ 和 $\Pi_h\boldsymbol\sigma$ 均有定义并满足这些等式.
有限元空间没有改变, 所以 inf--sup 条件也原样成立.

\MathBookSourcePage{266}
\begin{Lemma}\label{lem:incompressible-hhj-low-regularity-identities}
令 $\Omega$ 为有界平面多边形, 且 Stokes 问题的解
$(\mathbf u=\operatorname{curl}\psi,p)$ 满足
\begin{equation}\label{eq:incompressible-hhj-low-regularity-assumption}
\psi\in W^{3,r}(\Omega),
\qquad p\in W^{1,r}(\Omega),
\qquad r\in(1,2].
\end{equation}
则上述正交关系成立. 若此外 $\Omega$ 为凸域且
$\mathcal T_h$ 一致正则, 则前述稳定与误差估计也成立.
\end{Lemma}

\begin{Lemma}\label{lem:incompressible-hhj-low-regularity-restriction}
令 $\mathcal T_h$ 是 $\Omega$ 的正则剖分族. 对每个
$1<r\leq2$ 和每个对称张量
$\boldsymbol\tau\in W^{1,r}(\Omega)^4$, 有
\begin{equation}\label{eq:incompressible-hhj-low-regularity-restriction}
\|\Pi_h\boldsymbol\tau-\boldsymbol\tau\|_{0,\Omega}
\leq Ch^{2(1-1/r)}|\boldsymbol\tau|_{1,r,\Omega}.
\end{equation}
\end{Lemma}

\begin{Theorem}\label{thm:incompressible-hhj-low-regularity-error}
假设 \eqref{eq:incompressible-hhj-low-regularity-assumption} 成立.
若 $\mathcal T_h$ 正则, 则问题
\eqref{eq:incompressible-hhj-discrete-problem} 的解满足
\begin{equation}\label{eq:incompressible-hhj-low-regularity-tensor-error}
\|\boldsymbol\sigma-\boldsymbol\sigma_h\|_{0,\Omega}
\leq C_1h^{2(1-1/r)}|\psi|_{3,r,\Omega},
\end{equation}
其中
$\boldsymbol\sigma=(\partial^2\psi/\partial x_i\partial x_j)_{i,j}$.
若此外 $\mathcal T_h$ 一致正则且 $\Omega$ 为凸域,
则对每个实数 $\beta\geq2$, 存在 $C_2(\beta)>0$, 使
\begin{equation}\label{eq:incompressible-hhj-low-regularity-w1beta-error}
|\psi-\psi_h|_{1,\beta,\Omega}
\leq C_2(\beta)h^{2(1-1/r)}|\psi|_{3,r,\Omega}.
\end{equation}
当多项式次数 $l\geq2$ 时, 还有
\begin{equation}\label{eq:incompressible-hhj-low-regularity-piecewise-h2-error}
\left(\sum_{\kappa\in\mathcal T_h}
|\psi-\psi_h|_{2,\kappa}^2\right)^{1/2}
\leq C_3h^{2(1-1/r)}|\psi|_{3,r,\Omega}.
\end{equation}
\end{Theorem}

\subsection{压力的间断逼近}
\label{subsec:incompressible-hhj-discontinuous-pressure}

本小节分析从 Hellan--Herrmann--Johnson 格式恢复压力的方法.
压力通过问题 \eqref{eq:incompressible-hhj-velocity-pressure-weak}
的适当离散化得到; 当检验函数无散时, 它退化为
\eqref{eq:incompressible-hhj-discrete-problem}.
构造 $H_0(\operatorname{div};\Omega)$ 与 $L_0^2(\Omega)$
的有限维子空间 $D_{0h}$ 和 $Q_h$, 使离散无散向量恰为
$\operatorname{curl}\phi_h$, $\phi_h\in\Phi_h$, 并满足适当的 inf--sup 条件.

\MathBookSourcePage{267}
在参考单元 $\widehat\kappa$ 上, 若它是单位三角形, 取
\begin{equation}\label{eq:incompressible-hhj-pressure-reference-triangle}
D=P_{l-1}^2\oplus
\{p(\widehat{\mathbf x})\widehat{\mathbf x};
\ p\in\widetilde P_{l-1}\},
\end{equation}
其中 $\widetilde P_k$ 是 $k$ 次齐次多项式空间.
若 $\widehat\kappa$ 是单位正方形, 则取
\begin{equation}\label{eq:incompressible-hhj-pressure-reference-square}
D=Q_{l,l-1}\times Q_{l-1,l}.
\end{equation}
于是
\[
\begin{array}{lll}
\operatorname{div}D=P_{l-1},&
\ker(\operatorname{div})=\operatorname{curl}P_l,
&\widehat\kappa\text{ 为单位三角形},\\
\operatorname{div}D=Q_{l-1},&
\ker(\operatorname{div})=\operatorname{curl}Q_l,
&\widehat\kappa\text{ 为单位正方形}.
\end{array}
\]
可用下列自由度唯一确定 $\widehat{\mathbf v}\in D$:
\begin{equation}\label{eq:incompressible-hhj-pressure-reference-dofs}
\left\{
\begin{aligned}
&\int_{\widehat\kappa'}\widehat{\mathbf v}\cdot\widehat{\mathbf n}\,q\,ds
&&\forall q\in P_{l-1},
&&\widehat\kappa'\subset\partial\widehat\kappa,\\
&\int_{\widehat\kappa}\widehat{\mathbf v}\cdot\mathbf q\,dx
&&\forall\mathbf q\in
\begin{cases}
P_{l-2}^2,&\widehat\kappa\text{ 为三角形},\\
Q_{l-2,l-1}\times Q_{l-1,l-2},&\widehat\kappa\text{ 为正方形}.
\end{cases}
\end{aligned}
\right.
\end{equation}
每条边上的 $l$ 个自由度唯一决定该边上的法向分量.

由参考单元到任意单元 $\kappa$ 采用逆变变换
\begin{equation}\label{eq:incompressible-hhj-pressure-contravariant-transform}
\mathbf v\circ F_\kappa
=\frac1{J_\kappa}B_\kappa\widehat{\mathbf v},
\qquad
\widehat{\mathbf v}\circ F_\kappa^{-1}
=J_\kappa B_\kappa^{-1}\mathbf v.
\end{equation}
它保持散度、旋度和法向分量的相应积分结构.
在 $\Gamma_h$ 每条边上固定 $l$ 个不同点, 定义
\begin{equation}\label{eq:incompressible-hhj-pressure-velocity-space}
\begin{split}
D_{0h}=\{\mathbf v_h\in L^2(\Omega)^2;&
\ \mathbf v_h|_\kappa=F_\kappa\widehat{\mathbf v},
\ \widehat{\mathbf v}\in D;\\
&\mathbf v_h\cdot\mathbf n
\text{ 在每条内边的这些点连续, 并在边界点为零}\}.
\end{split}
\end{equation}

\MathBookSourcePage{268}
压力空间取为
\begin{equation}\label{eq:incompressible-hhj-pressure-space}
Q_h=\{q_h\in L_0^2(\Omega);
\ q_h|_\kappa\in P_{l-1}
\text{ 或 }Q_{l-1},
\text{ 依 }\kappa\text{ 为三角形或四边形而定}\}.
\end{equation}
立即有 $D_{0h}\subset H_0(\operatorname{div};\Omega)$, 且
\[
\{\mathbf v_h\in D_{0h};
\ (q_h,\operatorname{div}\mathbf v_h)=0\ \forall q_h\in Q_h\}
=\{\operatorname{curl}\phi_h;\ \phi_h\in\Phi_h\}.
\]

\begin{Remark}\label{rem:incompressible-hhj-pressure-lowest-order}
当 $l=1$ 时, $Q_h$ 的函数在各单元上为常数, 而参考空间函数为
$(c_1+c_3x_1,c_2+c_3x_2)$ (三角形) 或
$(c_1+c_3x_2,c_2+c_4x_1)$ (正方形).
\end{Remark}

由 \eqref{eq:incompressible-hhj-pressure-reference-dofs}
定义局部限制算子 $\widehat\pi$, 再由逆变变换得到
\begin{equation}\label{eq:incompressible-hhj-pressure-global-restriction}
(\pi_h\mathbf v)|_\kappa
=F_\kappa(\widehat\pi\widehat{\mathbf v})
\qquad\forall\kappa\in\mathcal T_h.
\end{equation}

\begin{Theorem}\label{thm:incompressible-hhj-pressure-restriction}
若 $\mathcal T_h$ 正则, 则对 $s\geq2$ 和
$\mathbf v\in H^1(\Omega)^2$,
\begin{equation}\label{eq:incompressible-hhj-pressure-restriction-w1s}
\left(\sum_{\kappa\in\mathcal T_h}
|\pi_h\mathbf v-\mathbf v|_{1,\kappa}^2\right)^{1/2}
+h^{-1}\|\pi_h\mathbf v-\mathbf v\|_{0,s,\Omega}
\leq C(s)|\mathbf v|_{1,s,\Omega}.
\end{equation}
此外,
\begin{equation}\label{eq:incompressible-hhj-pressure-restriction-divergence}
\|\operatorname{div}(\pi_h\mathbf v-\mathbf v)\|_{0,\Omega}
\leq C\|\operatorname{div}\mathbf v\|_{0,\Omega},
\end{equation}
并且对任意剖分,
\begin{equation}\label{eq:incompressible-hhj-pressure-restriction-orthogonality}
(q_h,\operatorname{div}(\pi_h\mathbf v-\mathbf v))=0
\qquad\forall q_h\in Q_h.
\end{equation}
若剖分仅由三角形组成, 则
\eqref{eq:incompressible-hhj-pressure-restriction-divergence}
以 $C=1$ 成立且不需正则性假设.
\end{Theorem}

\begin{Proof}
\eqref{eq:incompressible-hhj-pressure-restriction-divergence}
与 \eqref{eq:incompressible-hhj-pressure-restriction-orthogonality}
直接来自 $\pi_h$ 的定义. 为说明
\eqref{eq:incompressible-hhj-pressure-restriction-w1s},
考虑四边形单元, 令
$\widehat{\mathbf v}=F_\kappa^{-1}\mathbf v$ 并分解为
\begin{equation}\label{eq:incompressible-hhj-pressure-reference-decomposition}
\widehat{\mathbf v}=\operatorname{grad}\widehat q
+\operatorname{curl}\widehat\phi,
\end{equation}
其中
$\Delta\widehat q=\operatorname{div}\widehat{\mathbf v}$ 于
$\widehat\kappa$ 内, 且
$\widehat q|_{\partial\widehat\kappa}=0$.
参考正方形为凸域, 所以
$\|\widehat q\|_{2,\widehat\kappa}
\leq C\|\operatorname{div}\widehat{\mathbf v}\|_{0,\widehat\kappa}$,
并可选择 $\widehat\phi\in H^2(\widehat\kappa)$ 使
\[
\left\|\frac{\partial^2\widehat\phi}{\partial\widehat x_1^2}\right\|_{0,\widehat\kappa}^2
+\left\|\frac{\partial^2\widehat\phi}{\partial\widehat x_2^2}\right\|_{0,\widehat\kappa}^2
\leq C\left\{
\left\|\frac{\partial\widehat v_1}{\partial\widehat x_2}\right\|_{0,\widehat\kappa}^2
+\left\|\frac{\partial\widehat v_2}{\partial\widehat x_1}\right\|_{0,\widehat\kappa}^2
+\|\operatorname{div}\widehat{\mathbf v}\|_{0,\widehat\kappa}^2
\right\}.
\]
由附录逼近定理,
\begin{equation}\label{eq:incompressible-hhj-pressure-reference-estimate}
|\widehat\pi\widehat{\mathbf v}-\widehat{\mathbf v}|_{1,\widehat\kappa}
\leq C\left\{
\left\|\frac{\partial\widehat v_1}{\partial\widehat x_2}\right\|_{0,\widehat\kappa}^2
+\left\|\frac{\partial\widehat v_2}{\partial\widehat x_1}\right\|_{0,\widehat\kappa}^2
+\|\operatorname{div}\widehat{\mathbf v}\|_{0,\widehat\kappa}^2
\right\}^{1/2}.
\end{equation}
变换回 $\kappa$ 并使用剖分正则性即得结论.
\end{Proof}

\MathBookSourcePage{269}
Theorem~\ref{thm:incompressible-hhj-pressure-restriction}
蕴含下列 inf--sup 条件: 对每个 $q_h\in Q_h$, 存在
$\mathbf v_h\in D_{0h}$, 使
\[
(q_h,\operatorname{div}\mathbf v_h)=\|q_h\|_{0,\Omega}^2
\]
且对每个 $s\geq2$,
\[
\left(\sum_{\kappa\in\mathcal T_h}|\mathbf v_h|_{1,\kappa}^2\right)^{1/2}
+\|\mathbf v_h\|_{0,s,\Omega}
+\|\operatorname{div}\mathbf v_h\|_{0,\Omega}
\leq C(s)\|q_h\|_{0,\Omega}.
\]
用 $D_{0h}$ 和 $Q_h$ 离散化压力问题如下: 求 $p_h\in Q_h$, 使
\begin{equation}\label{eq:incompressible-hhj-discontinuous-pressure-problem}
\begin{split}
\int_\Omega p_h\operatorname{div}\mathbf v_h\,dx
={}&-\int_\Omega\mathbf f\cdot\mathbf v_h\,dx\\
&+\nu\sum_{\kappa\in\mathcal T_h}
\left\{\int_\kappa(\boldsymbol\lambda_h)_{ij}
\frac{\partial v_{hi}}{\partial x_j}\,dx
-\int_{\partial\kappa}M_{nt}(\boldsymbol\lambda_h)
\mathbf v_h\cdot\mathbf t\,ds\right\}
\end{split}
\end{equation}
对所有 $\mathbf v_h\in D_{0h}$ 成立, 其中
$\boldsymbol\lambda_h$ 由张量 $\boldsymbol\sigma_h$ 按
\eqref{eq:incompressible-hhj-tensor-correspondence} 确定.
由于 $\boldsymbol\sigma_h$ 可独立求出, 上述 inf--sup 条件保证
\eqref{eq:incompressible-hhj-discontinuous-pressure-problem}
在 $Q_h$ 中有唯一解.

\MathBookSourcePage{270}
\begin{Theorem}\label{thm:incompressible-hhj-discontinuous-pressure-error}
设 $\Omega$ 为有界多边形, $\mathcal T_h$ 是其一致正则剖分.
若 Stokes 问题的解满足
\[
\mathbf u\in H^{k+1}(\Omega)^2,
\qquad p\in H^k(\Omega)\cap L_0^2(\Omega),
\qquad k\in[1,l],
\]
则问题 \eqref{eq:incompressible-hhj-discontinuous-pressure-problem}
的解满足
\begin{equation}\label{eq:incompressible-hhj-discontinuous-pressure-error}
\|p-p_h\|_{0,\Omega}
\leq Ch^k\{\|\mathbf u\|_{k+1,\Omega}+|p|_{k,\Omega}\}.
\end{equation}
\end{Theorem}

\begin{Proof}
对每个 $q_h\in Q_h$, 利用
\eqref{eq:incompressible-hhj-tensor-correspondence} 得
\[
\begin{split}
\left|\int_\Omega(q_h-p_h)\operatorname{div}\mathbf v_h\,dx\right|
\leq{}&\|q_h-p\|_{0,\Omega}\|\operatorname{div}\mathbf v_h\|_{0,\Omega}\\
&+\nu\|\boldsymbol\sigma-\boldsymbol\sigma_h\|_{0,\Omega}
\left(\sum_{\kappa\in\mathcal T_h}|\mathbf v_h|_{1,\kappa}^2\right)^{1/2}\\
&+\nu\|M_n(\boldsymbol\sigma-\boldsymbol\sigma_h)\|_{0,\Gamma_h}
\|S(\mathbf v_h\cdot\mathbf t)\|_{0,\Gamma_h}.
\end{split}
\]
由 inf--sup 条件选择 $\mathbf v_h\in D_{0h}$, 得
\[
\|q_h-p_h\|_{0,\Omega}
\leq C\{\|q_h-p\|_{0,\Omega}
+\nu\|\boldsymbol\sigma-\boldsymbol\sigma_h\|_{0,h}\}.
\]
这里写 $\mathbf v_h=\pi_h\mathbf v$, 其中
$\operatorname{div}\mathbf v=q_h-p_h$,
$|\mathbf v|_{1,\Omega}\leq C\|q_h-p_h\|_{0,\Omega}$,
$\mathbf v\in H_0^1(\Omega)^2$, 并利用
$S(\mathbf v\cdot\mathbf t)=0$ 与
\eqref{eq:incompressible-hhj-pressure-reference-estimate}
估计边界项. 最后应用张量误差估计和附录的压力逼近估计,
得到 \eqref{eq:incompressible-hhj-discontinuous-pressure-error}.
注意, 只在张量误差估计中需要剖分的一致正则性.
\end{Proof}

\begin{Remark}\label{rem:incompressible-hhj-pressure-with-vorticity-scheme}
该压力间断逼近也可与 Section~\ref{sec:incompressible-sfvp}
的 stream function--vorticity 格式结合. 此时求由
\eqref{eq:incompressible-hhj-pressure-space} 定义的 $p_h\in Q_h$, 使
\begin{equation}\label{eq:incompressible-vorticity-discontinuous-pressure}
(p_h,\operatorname{div}\mathbf v_h)
=\nu(\operatorname{curl}\omega_h,\mathbf v_h)
-(\mathbf f,\mathbf v_h)
\qquad\forall\mathbf v_h\in D_{0h}.
\end{equation}
该格式与前述连续压力逼近具有相同收敛阶,
但证明要求解具有更高正则性, 参见 Girault \& Raviart~[33].
\end{Remark}

\MathBookSourcePage{271}
\begin{Remark}\label{rem:incompressible-hhj-continuous-pressure-alternative}
也可采用类似于前述连续压力格式的逼近,
但相应格式的误差分析更为精细.
\end{Remark}
